\documentclass[a4paper,12pt]{article}
\usepackage[utf8]{inputenc}
\usepackage[T1]{fontenc}
\usepackage[french]{babel}
\usepackage{amsmath,amsfonts,amsthm,amssymb,amsbsy,amscd,mathrsfs}
\usepackage{lmodern}
\usepackage{fourier}
%\usepackage{times,mathptmx}
%\DeclareSymbolFont{lmodern}{U}{euex}{m}{n}
%\DeclareMathSymbol\infty \mathord{lmodern}{"31}
\usepackage[scaled=0.875]{helvet} % ss
\usepackage{textcomp}
\usepackage{dsfont} % pour les symboles des ensembles
\usepackage[left=2cm,right=2cm,top=1.5cm,bottom=2.5cm]{geometry}
\usepackage{fancyhdr} 
\fancyhf{}
\setlength{\headheight}{26pt}
\setlength{\footskip}{18pt}
\renewcommand {\headrulewidth}{0pt}
\renewcommand {\footrulewidth}{0pt}
%\usepackage{lastpage}
%\usepackage{makeidx}
\usepackage{graphicx} % inclure le fichier graphique de pondichery 2010
\usepackage{hyperxmp}
\usepackage[colorlinks=true,pdfstartview=FitV,linkcolor=blue,citecolor=blue,urlcolor=blue]{hyperref}
\pagestyle{fancy}
\usepackage{array}
\usepackage{tabularx}
\usepackage{multirow}
\usepackage{hhline}
\usepackage[table]{xcolor}
\usepackage{arydshln}%%% lignes en tirets dans un tableau 
\usepackage{mathtools} %pour l'alignement des termes d'une matrice
\usepackage{fltpoint}
\usepackage{xcolor,colortbl}
\usepackage{pstricks,pst-all,pst-plot,pstricks-add,pst-tree,pst-3dplot,pst-func}
\usepackage{ifthen}
\usepackage{calc}
\usepackage{esvect}
\usepackage[np]{numprint}
\setlength{\parindent}{0mm}
\usepackage [alwaysadjust]{paralist}
\setdefaultenum {1.}{a)}{}{}
\hyphenpenalty=10000 
\tolerance=2000
\emergencystretch=1em
\widowpenalty=10000
\clubpenalty=10000
\raggedbottom
\renewcommand*{\tabularxcolumn}[1]{m{#1}} %centrage vertical des cellules d'un tableau tabularx
\newcommand{\Oij}{$\left( {{\mathrm{O}};\vec i,\vec j} \right)$}
\newcommand{\Oijk}{$\left( {{\mathrm{O}};\vec i,\vec j ,\vec k} \right)$}
\newcommand{\euro}{\texteuro{}}
\newcommand{\R}{\mathds {R}}
%\newcommand{\C}{\mathds {C}}
\newcommand{\N}{\mathds {N}}
\newcommand{\Z}{\mathds {Z}}
\newcommand{\Q}{\mathds {Q}}
\newcommand{\e}{\mathrm {e}}
\newcommand{\dd}{\,\mathrm{d}}
\DecimalMathComma
%%% Mise en forme des algorithmes
\usepackage[french,boxed,titlenumbered,lined,longend]{algorithm2e}
  \SetKwIF {Si}{SinonSi}{Sinon}{si}{alors}{sinon\_si}{alors}{fin~si}
 \SetKwFor{Tq}{tant\_que~}{~faire~}{fin~tant\_que}
 \SetKwFor{PourCh}{pour\_chaque }{ faire }{fin pour\_chaque}
 \SetKwInput{Sortie}{Sortie}
\newcommand{\Algocmd}[1]{\textsf{\textsc{\textbf{#1}}}}\SetKwSty{Algocmd}
  \newcommand{\AlgCommentaire}[1]{\textsl{\small  #1}}  
%%%%%%%%%%%%%%%
\newcounter{nexo}    % déclaration du numéro d'exo
\setcounter{nexo}{0} % initialisation du numero
\newcommand{\exo}{%
  \stepcounter{nexo} 
  \par{\textsf{\textbf{\textcolor{blue}{{\textsc{exercice  \arabic{nexo}}}}}}}
  }%
\newcommand{\QCM}[3]{ $\bullet$ \quad #1 \hfill 
$\bullet$ \quad #2 \hfill 
$\bullet$ \quad #3 }
\newcommand{\QCMc}[4]{   #1 \\ 
\makebox[3cm][l]{\psframe(.25,.25) \qquad #2 }\hfill 
\makebox[3cm][l]{\psframe(.25,.25) \qquad #3 }\hfill 
\makebox[3cm][l]{\psframe(.25,.25) \qquad #4 }} % QCM avec cases à cocher%
% en-tête
%\lhead{\footnotesize{\textsf{Lycée JANSON DE SAILLY\\ {\today}}}} %gauche
\chead{\textsf{\textcolor{red}{\textbf{\textsc{DS01}}}}} %centre
%\rhead{\footnotesize{\fbox{\textsf{S1E2}}}} %droit
%pied de page
%\lfoot{\scriptsize{\textsf{\textsc{A. Yallouz}}}} %gauche
\cfoot{\scriptsize{\textsf{-~\thepage{}~-}}}  %centre                                                   
%\rfoot{\htmladdnormallink{\scriptsize\textsf{\textsc{MATH@ES}}}{http://yallouz.arie.free.fr}} %droit 
\hypersetup{
pdftitle={titre du document.}, 
pdfauthor={Arié Yallouz},   
pdfsubject={Exercices Terminale ES.}, 
pdfkeywords={exercices, sujets.},
pdfcontacturl={http//yallouz.arie.free.fr/},
pdflang={fr},
baseurl={http://yallouz.arie.free.fr/serveurtex/exostex.php}
}
\begin{document}

\exo

Soit l'équation différentielle \qquad $(E): \quad y'-3y=-21x+13$


\begin{enumerate}
	\item Montrer que la fonction $f_p$ définie sur $\R$ par $f_p(x)=7x-2$ est une solution particulière de $(E)$
	
	\vspace{6cm}
	\item Déterminer les solutions de l'équation homogène \qquad $(E_0): \quad y'-3y=0$
	
		\vspace{4cm}

	\item En déduire que les solutions de $(E)$ peuvent s'écrire sous la forme $$f(x)=7x-2+k\e^{3x}$$
	
		\vspace{4cm}

	
	\item Déterminer la valeur de la constante $k$ telle que $f(0)=5$
\end{enumerate}

\newpage

\exo 

\medskip


On se propose de résoudre l'équation différentielle $$(E): \quad y'+y=\e^{-x}$$

\begin{enumerate}
	
	\item Résoudre dans $\R$ l'équation homogène $(E_0): \quad y'+y=0$
	
	\vspace{3cm}
	
	\item Montrer que la fonction $f_p$ définie sur $\R$ par $$f_p(x)=x\e^{-x}$$ est une solution particulière de $(E)$
	
	\vspace{6cm}
	
	\item En déduire la solution générale $f$ de $(E)$
	
	\vspace{2cm}

\item Déterminer la solution particulière de $(E)$ qui est telle que la courbe représentative de la fonction $f$ solution de $(E)$ passe par le point $(0;3)$
	\end{enumerate}







\newpage

\begin{center}
\begin{Huge}
CORRECTION DS01

\end{Huge}
\end{center}


\setcounter{nexo}{0}

\exo

\begin{enumerate}
	\item $f_p(x)=7x-2$ et $f'_p(x)=7$
	
	Donc $f'_p(x)-3f_p(x)=7-3(7x-2)=7-21x+6=-21x+13$
	
	Donc $f_p$ est bien une solution particulière de l'équation différentielle $(E)$.
	
	\bigskip
	
	\item Les solutions de l'équation homogène $(E_0)$ sont toutes les fonctions du type $$f_0(x)=k\e^{3x}, \quad k \in \R$$
	
	\bigskip
	
	\item Les solutions de l'équation $(E)$ sont donc les fonctions $f$ de la forme:
	
	$$f(x)=f_p(x)+f_0(x)=7x-2+k\e^{3x} \quad k\in\R$$
	
	\bigskip
	
	\item On cherche la valeur du réel $k$ tel que $f(0)=5$
	
	On a $f(0)=7 \times 0 - 2 + k \e^{3 \times 0} = -2 + k$ Donc finalement $k=7$
	
	$$f(x)=7x-2+5\e^{3x}$$
\end{enumerate}

\bigskip

\exo

\bigskip

\begin{enumerate}
	\item Les solutions de l'équation homogène $(E_0)$ sont toutes les fonctions du type $$f_0(x)=k\e^{-x}, \quad k \in \R$$

\bigskip


	\item $f_p(x)=x\e^{-x}$ et $f'_p(x)=(1-x)\e^{-x}$
	
	Donc $f'_p(x)+f_p(x)=(1-x)\e{-x}+x\e^{-x}=\e^{-x}-x\e^{-x}+x\e^{-x}=\e^{-x}$
	
	Donc $f_p$ est bien une solution particulière de l'équation différentielle $(E)$.
	
	\bigskip
	
	
	
	\item Les solutions de l'équation $(E)$ sont donc les fonctions $f$ de la forme:
	
	$$f(x)=f_p(x)+f_0(x)=x\e^{-x}+k\e^{-x}=(x-k)\e^{-x} \quad k\in\R$$
	
	\bigskip
	
	\item On cherche la valeur du réel $k$ tel que $f(0)=3$
	
	On a $f(0)=k$ Donc finalement $k=3$
	
	$$f(x)=(x-3)\e^{-x}$$
\end{enumerate}
\end{document}
